\documentclass[11pt]{article}
\usepackage[margin=1in]{geometry}
\usepackage{amsmath}
\usepackage{amssymb}
\usepackage{hyperref}
\usepackage{mdframed}
\usepackage{abstract}
\usepackage{titlesec}

\titleformat{\section}{\large\bfseries}{\thesection.}{0.5em}{}
\titleformat{\subsection}{\normalsize\bfseries}{\thesubsection.}{0.5em}{}

\title{\textbf{Intuition for Diffusion in Nature and in Generative Modelling}\\[0.4em]
\large Lecture Notes}
\author{Carlos Cotrini\thanks{Department of Computer Science, ETH Z\"{u}rich.}
\and
Juan Tarazona\thanks{Department of Mechanical Engineering, ETH Z\"{u}rich.}}
\date{}

\begin{document}
\maketitle

\begin{abstract}
These notes provide an accessible introduction to diffusion processes and their role in modern generative modelling.
We begin with physical intuition, develop the necessary probabilistic tools at an elementary level, and proceed to a self-contained derivation of the denoising diffusion probabilistic model (DDPM) framework.
Proofs are collected in dedicated boxes and may be skipped on a first reading without loss of continuity.
\end{abstract}

\tableofcontents
\bigskip

% -----------------------------------------------------------------------
\section{Physical Diffusion and the Ubiquity of Bell-Shaped Distributions}
% -----------------------------------------------------------------------

Many observable phenomena arise from the cumulative effect of a large number of small, independent random influences.
The following examples illustrate this principle across different physical and social contexts.

\begin{itemize}
  \item \textbf{Dye spreading in water.}
        Each dye molecule is displaced by countless collisions with surrounding water molecules.
        After many such collisions, the dye cloud evolves into a smooth, expanding blur whose profile is well described by a bell-shaped distribution.
  \item \textbf{Heat conduction in a metal rod.}
        Thermal energy propagates through a lattice via many microscopic interactions, producing a smooth temperature profile that spreads over time.
  \item \textbf{Measurement noise.}
        When a sensor is subject to many small, independent error sources, the aggregate error tends to be bell-shaped around zero.
  \item \textbf{Human height.}
        Adult height is influenced by many genetic and environmental factors.
        When numerous small additive effects combine, the resulting population distribution is approximately bell-shaped around a central value.
\end{itemize}

\noindent
The common thread is the following principle: \emph{when an outcome is the sum of many small, independent pushes, the distribution of that outcome tends to be bell-shaped}.
We will call this bell-shaped family the \textbf{normal curve} (in one dimension) and the \textbf{normal cloud} (in multiple dimensions).

% -----------------------------------------------------------------------
\section{Probabilities as Integrals: the One-Dimensional Case}
% -----------------------------------------------------------------------

Suppose a scalar quantity $x$ follows a normal curve centred at $\mu$ with \textbf{spread} $\sigma$ (``sigma'').
The spread $\sigma$ governs the width of the bell: a larger value of $\sigma$ corresponds to outcomes that are more dispersed around the centre.
As a rule of thumb, the interval $[\mu - 2\sigma,\, \mu + 2\sigma]$ captures approximately $95\%$ of all outcomes.

In $d$ dimensions, the analogous statement is that a sphere of radius
\[
  r_{95} \;\approx\; \sigma + \sqrt{d}
\]
centred at $\mu$ captures approximately $95\%$ of the cloud.
The additional $\sqrt{d}$ term reflects the geometric fact that, in high dimensions, the probability mass of an isotropic cloud concentrates on a thin shell rather than near the centre.

\subsection*{Probability of landing in an interval}

The probability that $x$ falls between $a$ and $b$ is given by the area under the normal curve over that interval:
\[
  \Pr(a \le x \le b)
  \;=\;
  \int_{a}^{b} p(x)\,dx,
\]
where $p(x)$ denotes the height of the normal curve at the point $x$.
Geometrically, the integral accumulates the areas of infinitesimally thin rectangles of height $p(x)$ and width $dx$.
In practice, this integral is evaluated numerically or via lookup tables, but the conceptual content is simply \emph{area under the bell}.

% -----------------------------------------------------------------------
\section{Probabilities as Integrals: the Multidimensional Case}
% -----------------------------------------------------------------------

Let $x$ now be a vector in $\mathbb{R}^d$ (for example, a two-dimensional spatial location, or the pixel values of an image).
The normal curve generalises to a \emph{normal cloud}, which assigns a probability density $p(x)$ to each point in $\mathbb{R}^d$.

\subsection*{Probability of landing in a region}

For any region $R \subset \mathbb{R}^d$, the probability that $x$ falls in $R$ is
\[
  \Pr(x \in R)
  \;=\;
  \int_{x \in R} p(x)\,dx.
\]
The one-dimensional notion of ``area under a curve'' is replaced by ``volume under a surface'' (in two dimensions) or, more generally, by the integral of probability mass over the region $R$.

\subsection*{A canonical special case: probability mass inside a ball}

When the cloud is isotropic (equally spread in every direction) and centred at the origin, a natural quantity is the fraction of mass contained within a ball of radius $r$:
\[
  \Pr(\|x\| \le r) \;=\; \int_{\|x\| \le r} p(x)\,dx.
\]
This quantity plays a central role in understanding the geometry of high-dimensional distributions.

% -----------------------------------------------------------------------
\section{Simulating Bell-Shaped Noise from Uniform Draws}
% -----------------------------------------------------------------------

To simulate outcomes from bell-shaped distributions, we require a practical method to generate the necessary random noise.
We assume access only to a \textbf{uniform random number generator} producing values in $[-1,1]$.

\subsection*{The summation method}

A classical observation is the following:
\begin{quote}
  \emph{The sum of many independent uniform random variables increasingly resembles a bell-shaped distribution as the number of terms grows.}
\end{quote}

Let $u_1, u_2, \ldots, u_n$ be independent draws from $\mathrm{Uniform}[-1,1]$, each with mean $0$.
Mathematics guarantees that the rescaled sum
\[
  z \;=\; (u_1 + \cdots + u_n) \cdot \sqrt{\tfrac{3}{n}}
\]
behaves as a standard bell-shaped noise for sufficiently large $n$.
A practical choice is $n = 12$, which simplifies the formula to $z = (u_1 + \cdots + u_{12})/2$.

\noindent
To obtain a bell-shaped value centred at $\mu$ with spread $\sigma$, one then computes
\[
  x = \mu + \sigma z.
\]

\subsection*{Multidimensional sampling: the isotropic cloud}

The simplest multivariate normal cloud is the \textbf{isotropic} case, in which the distribution is equally spread in every direction and coordinates are mutually independent.
Sampling from such a cloud of dimension $d$ proceeds as follows:
\begin{enumerate}
  \item Generate $d$ independent standard bell-shaped values $z_1, \ldots, z_d$ using the summation method above.
  \item Assemble them into a vector $z = (z_1, \ldots, z_d)^\top$.
  \item Apply the affine transformation $x = \mu + \sigma z$.
\end{enumerate}
The result is a spherical cloud centred at $\mu$, with each coordinate independently dispersed with spread $\sigma$.

% -----------------------------------------------------------------------
\section{The Diffusion Process: Forward Corruption and Reverse Restoration}
% -----------------------------------------------------------------------

A \textbf{diffusion process}, in the sense used in generative modelling, is a pair of stochastic procedures:
\begin{itemize}
  \item \textbf{Forward process:} structured data $x_0$ is gradually corrupted by the successive injection of small amounts of noise, producing a sequence $x_1, x_2, \ldots, x_T$ that converges to pure noise.
  \item \textbf{Reverse process:} a learned model progressively removes noise from a noisy sample, recovering structured data from pure noise.
\end{itemize}

Throughout, we restrict attention to \textbf{sampling equations}: explicit recipes for generating the next state given the current one.

\subsection*{Schedule and notation}

We fix a \textbf{noise schedule} $\beta_1, \beta_2, \ldots, \beta_T$ with each $\beta_t \in (0,1)$ small, and define
\[
  \alpha_t \;=\; 1 - \beta_t, \qquad
  \bar{\alpha}_t \;=\; \prod_{s=1}^{t} \alpha_s.
\]
The quantity $\bar{\alpha}_t$ represents the cumulative retention of the original signal after $t$ steps.
We also fix a source of \textbf{standard noise} $\varepsilon$: in one dimension a scalar, in $d$ dimensions a vector with independent standard bell-shaped coordinates.

% -----------------------------------------------------------------------
\section{The Forward Process}
% -----------------------------------------------------------------------

\subsection*{One-step forward transition}

Given the state $x_{t-1}$, the next corrupted state $x_t$ is obtained by
\[
  x_t \;=\; \sqrt{\alpha_t}\,x_{t-1} \;+\; \sqrt{1-\alpha_t}\,\varepsilon_t,
\]
where $\varepsilon_t$ is a freshly sampled standard noise vector.
The factor $\sqrt{\alpha_t}$ attenuates the signal, while the term $\sqrt{1-\alpha_t}\,\varepsilon_t$ injects a controlled amount of fresh randomness.

\subsection*{Multi-step forward transition}

Applying the one-step rule repeatedly, one may express $x_t$ directly in terms of the original $x_0$:
\[
  x_t \;=\; \sqrt{\bar{\alpha}_t}\,x_0 \;+\; \sqrt{1-\bar{\alpha}_t}\,\varepsilon,
\]
where $\varepsilon$ is a single standard noise vector.
As $t$ increases, $\bar{\alpha}_t \to 0$ and $x_t$ converges to pure noise.

\begin{mdframed}[linewidth=1pt, frametitle={\textbf{Proof (for the curious reader): multi-step forward formula}}]
\textit{(This derivation may be skipped on a first reading.)}

\medskip
\noindent
\textbf{Claim.} Applying the one-step rule $t$ times yields
$x_t = \sqrt{\bar\alpha_t}\,x_0 + \sqrt{1-\bar\alpha_t}\,\varepsilon$,
where $\bar\alpha_t = \prod_{s=1}^{t}\alpha_s$ and $\varepsilon$ is standard noise.

\medskip\noindent
\emph{Base case ($t=1$).}
By the one-step rule, $x_1 = \sqrt{\alpha_1}\,x_0 + \sqrt{1-\alpha_1}\,\varepsilon_1$, which matches the claim with $\bar\alpha_1 = \alpha_1$.

\medskip\noindent
\emph{Inductive step.}
Assume the claim holds at step $t-1$:
$x_{t-1} = \sqrt{\bar\alpha_{t-1}}\,x_0 + \sqrt{1-\bar\alpha_{t-1}}\,\varepsilon'$.
Applying the one-step rule:
\[
x_t
= \sqrt{\alpha_t}\!\left(\sqrt{\bar\alpha_{t-1}}\,x_0 + \sqrt{1-\bar\alpha_{t-1}}\,\varepsilon'\right) + \sqrt{1-\alpha_t}\,\varepsilon_t
= \sqrt{\bar\alpha_t}\,x_0
  + \underbrace{\sqrt{\alpha_t(1-\bar\alpha_{t-1})}\,\varepsilon' + \sqrt{1-\alpha_t}\,\varepsilon_t}_{\text{sum of two independent bell-shaped noises}}.
\]
The sum of two independent bell-shaped noises with variances $a^2$ and $b^2$ is a bell-shaped noise with variance $a^2+b^2$.
Here the combined variance is $\alpha_t(1-\bar\alpha_{t-1}) + (1-\alpha_t) = 1 - \bar\alpha_t$,
so $x_t = \sqrt{\bar\alpha_t}\,x_0 + \sqrt{1-\bar\alpha_t}\,\varepsilon$. \hfill$\square$
\end{mdframed}

% -----------------------------------------------------------------------
\section{The Reverse Process}
% -----------------------------------------------------------------------

\subsection*{Step 1: the ideal reverse step (assuming $\varepsilon$ is known)}

Recall the multi-step forward formula:
\[
  x_t \;=\; \sqrt{\bar{\alpha}_t}\,x_0 \;+\; \sqrt{1-\bar{\alpha}_t}\,\varepsilon.
\]
Were the noise $\varepsilon$ known, one could recover $x_0$ exactly:
\[
  x_0 \;=\; \frac{x_t - \sqrt{1-\bar{\alpha}_t}\,\varepsilon}{\sqrt{\bar{\alpha}_t}}.
\]
Substituting this expression into the forward formula at step $t-1$ and taking the conditional mean yields the \emph{ideal} one-step reverse formula:
\[
  x_{t-1}
  \;=\;
  \frac{1}{\sqrt{\alpha_t}}
  \!\left(
    x_t \;-\; \frac{\beta_t}{\sqrt{1-\bar{\alpha}_t}}\,\varepsilon
  \right)
  \;+\; \sqrt{\beta_t}\,z,
\]
where $z$ is a fresh standard noise injection that maintains the correct stochastic dynamics.
This formula is exact when $\varepsilon$ is available.

\subsection*{Step 2: the fundamental difficulty at inference time}

During generation, one begins from pure noise and has no access to a real datum $x_0$.
Consequently, the noise vector $\varepsilon$ that would have produced the observed $x_t$ is unknown and must be estimated.

\subsection*{Step 3: learned denoising via $\widehat{\varepsilon}_\theta$}

A neural network $\widehat{\varepsilon}_\theta(x_t, t)$ is trained to estimate $\varepsilon$ given only $x_t$ and the time index $t$.
Substituting this estimate in place of the unknown $\varepsilon$ yields the practical reverse step:
\[
  x_{t-1}
  \;=\;
  \frac{1}{\sqrt{\alpha_t}}
  \!\left(
    x_t \;-\; \frac{\beta_t}{\sqrt{1-\bar{\alpha}_t}} \cdot \widehat{\varepsilon}_\theta(x_t,t)
  \right)
  \;+\;
  \sqrt{\beta_t}\, z,
\]
where:
\begin{itemize}
  \item $\widehat{\varepsilon}_\theta(x_t,t)$ is the network's estimate of the noise $\varepsilon$ present in $x_t$;
  \item $\dfrac{\beta_t}{\sqrt{1-\bar{\alpha}_t}}$ is a known scalar determined entirely by the schedule;
  \item $z$ is fresh standard noise, set to zero on the final step $t=1$;
  \item $\sqrt{\beta_t}$ governs the magnitude of the stochastic component of the reverse step.
\end{itemize}

\begin{mdframed}[linewidth=1pt, frametitle={\textbf{Proof (for the curious reader): derivation of the reverse step}}]
\textit{(This derivation may be skipped on a first reading.)}

\medskip\noindent
\textbf{Goal.} Derive the reverse step formula from the multi-step forward formula and the noise schedule.

\medskip\noindent
\textbf{1. Recover $x_0$ from $x_t$ and $\varepsilon$.}
From $x_t = \sqrt{\bar\alpha_t}\,x_0 + \sqrt{1-\bar\alpha_t}\,\varepsilon$, solving for $x_0$:
\[
x_0 = \frac{x_t - \sqrt{1-\bar\alpha_t}\,\varepsilon}{\sqrt{\bar\alpha_t}}.
\tag{$*$}
\]

\medskip\noindent
\textbf{2. Express the conditional mean of $x_{t-1}$ given $x_t$.}
The multi-step formula at step $t-1$ gives $x_{t-1} = \sqrt{\bar\alpha_{t-1}}\,x_0 + \sqrt{1-\bar\alpha_{t-1}}\,\varepsilon'$.
Conditioning on $x_t$ (i.e.\ fixing $x_0$ via $(*)$ and suppressing the residual $\varepsilon'$):
\[
\mathbb{E}[x_{t-1} \mid x_t]
= \sqrt{\bar\alpha_{t-1}}\cdot\frac{x_t - \sqrt{1-\bar\alpha_t}\,\varepsilon}{\sqrt{\bar\alpha_t}}.
\]

\medskip\noindent
\textbf{3. Simplify using $\bar\alpha_t = \alpha_t\,\bar\alpha_{t-1}$.}
Since $\sqrt{\bar\alpha_{t-1}}/\sqrt{\bar\alpha_t} = 1/\sqrt{\alpha_t}$, one obtains
\[
\mathbb{E}[x_{t-1} \mid x_t]
= \frac{1}{\sqrt{\alpha_t}}\!\left(x_t - \frac{\beta_t}{\sqrt{1-\bar\alpha_t}}\,\varepsilon\right),
\]
where the coefficient of $\varepsilon$ follows from the identity $\frac{\sqrt{1-\bar\alpha_t}}{\sqrt{\alpha_t}} - \sqrt{\frac{1-\bar\alpha_{t-1}}{\alpha_t}} = \frac{\beta_t}{\sqrt{1-\bar\alpha_t}}$.

\medskip\noindent
\textbf{4. Restore the stochastic component.}
The posterior distribution of $x_{t-1}$ given $x_t$ and $x_0$ is a normal cloud centred at the mean above, with variance $\beta_t$ (the choice made in the original DDPM formulation).
Adding the corresponding noise term $\sqrt{\beta_t}\,z$, $z\sim\mathcal{N}(0,I)$, yields
\[
x_{t-1}
= \frac{1}{\sqrt{\alpha_t}}\!\left(x_t - \frac{\beta_t}{\sqrt{1-\bar\alpha_t}}\,\varepsilon\right) + \sqrt{\beta_t}\,z.
\]
Replacing $\varepsilon$ by the network estimate $\widehat\varepsilon_\theta(x_t,t)$ gives the formula used in practice. \hfill$\square$
\end{mdframed}

\subsection*{Multi-step generation procedure}

Starting from pure noise, the full generation algorithm proceeds as follows:
\begin{enumerate}
  \item Sample $x_T$ from a standard normal cloud.
  \item For $t = T, T-1, \ldots, 1$, compute
  \[
  x_{t-1}
  \;=\;
  \frac{1}{\sqrt{\alpha_t}}
  \!\left(
    x_t \;-\; \frac{\beta_t}{\sqrt{1-\bar{\alpha}_t}} \cdot \widehat{\varepsilon}_\theta(x_t,t)
  \right)
  \;+\;
  \sqrt{\beta_t}\, z.
  \]
  \item Return $x_0$ as the generated sample.
\end{enumerate}

% -----------------------------------------------------------------------
\section{Training the Diffusion Model}
% -----------------------------------------------------------------------

The network $\widehat{\varepsilon}_\theta$ is trained by minimising the mean squared error between its predictions and the true noise vectors used in the forward process.
Concretely, training proceeds by repeatedly:
\begin{itemize}
  \item drawing a real datum $x_0$ from the training set;
  \item sampling a time step $t$ uniformly from $\{1, \ldots, T\}$ and a noise vector $\varepsilon$;
  \item corrupting $x_0$ to obtain $x_t = \sqrt{\bar{\alpha}_t}\,x_0 + \sqrt{1-\bar{\alpha}_t}\,\varepsilon$;
  \item updating $\theta$ to minimise $\|\varepsilon - \widehat{\varepsilon}_\theta(x_t, t)\|^2$.
\end{itemize}

The key insight is that the training objective is entirely self-supervised: the targets $\varepsilon$ are generated synthetically by the forward process and require no human annotation.
The trained network encodes, for each noise level $t$, the mapping from a noisy observation $x_t$ back to the noise $\varepsilon$ that produced it --- which is precisely the information needed to run the reverse process.

\end{document}
